% Options for packages loaded elsewhere
\PassOptionsToPackage{unicode}{hyperref}
\PassOptionsToPackage{hyphens}{url}
\documentclass[
  british,
]{article}
\usepackage{xcolor}
\usepackage[margin=1in]{geometry}
\usepackage{amsmath,amssymb}
\setcounter{secnumdepth}{-\maxdimen} % remove section numbering
\usepackage{iftex}
\ifPDFTeX
  \usepackage[T1]{fontenc}
  \usepackage[utf8]{inputenc}
  \usepackage{textcomp} % provide euro and other symbols
\else % if luatex or xetex
  \usepackage{unicode-math} % this also loads fontspec
  \defaultfontfeatures{Scale=MatchLowercase}
  \defaultfontfeatures[\rmfamily]{Ligatures=TeX,Scale=1}
\fi
\usepackage{lmodern}
\ifPDFTeX\else
  % xetex/luatex font selection
\fi
% Use upquote if available, for straight quotes in verbatim environments
\IfFileExists{upquote.sty}{\usepackage{upquote}}{}
\IfFileExists{microtype.sty}{% use microtype if available
  \usepackage[]{microtype}
  \UseMicrotypeSet[protrusion]{basicmath} % disable protrusion for tt fonts
}{}
\makeatletter
\@ifundefined{KOMAClassName}{% if non-KOMA class
  \IfFileExists{parskip.sty}{%
    \usepackage{parskip}
  }{% else
    \setlength{\parindent}{0pt}
    \setlength{\parskip}{6pt plus 2pt minus 1pt}}
}{% if KOMA class
  \KOMAoptions{parskip=half}}
\makeatother
\usepackage{color}
\usepackage{fancyvrb}
\newcommand{\VerbBar}{|}
\newcommand{\VERB}{\Verb[commandchars=\\\{\}]}
\DefineVerbatimEnvironment{Highlighting}{Verbatim}{commandchars=\\\{\}}
% Add ',fontsize=\small' for more characters per line
\usepackage{framed}
\definecolor{shadecolor}{RGB}{248,248,248}
\newenvironment{Shaded}{\begin{snugshade}}{\end{snugshade}}
\newcommand{\AlertTok}[1]{\textcolor[rgb]{0.94,0.16,0.16}{#1}}
\newcommand{\AnnotationTok}[1]{\textcolor[rgb]{0.56,0.35,0.01}{\textbf{\textit{#1}}}}
\newcommand{\AttributeTok}[1]{\textcolor[rgb]{0.13,0.29,0.53}{#1}}
\newcommand{\BaseNTok}[1]{\textcolor[rgb]{0.00,0.00,0.81}{#1}}
\newcommand{\BuiltInTok}[1]{#1}
\newcommand{\CharTok}[1]{\textcolor[rgb]{0.31,0.60,0.02}{#1}}
\newcommand{\CommentTok}[1]{\textcolor[rgb]{0.56,0.35,0.01}{\textit{#1}}}
\newcommand{\CommentVarTok}[1]{\textcolor[rgb]{0.56,0.35,0.01}{\textbf{\textit{#1}}}}
\newcommand{\ConstantTok}[1]{\textcolor[rgb]{0.56,0.35,0.01}{#1}}
\newcommand{\ControlFlowTok}[1]{\textcolor[rgb]{0.13,0.29,0.53}{\textbf{#1}}}
\newcommand{\DataTypeTok}[1]{\textcolor[rgb]{0.13,0.29,0.53}{#1}}
\newcommand{\DecValTok}[1]{\textcolor[rgb]{0.00,0.00,0.81}{#1}}
\newcommand{\DocumentationTok}[1]{\textcolor[rgb]{0.56,0.35,0.01}{\textbf{\textit{#1}}}}
\newcommand{\ErrorTok}[1]{\textcolor[rgb]{0.64,0.00,0.00}{\textbf{#1}}}
\newcommand{\ExtensionTok}[1]{#1}
\newcommand{\FloatTok}[1]{\textcolor[rgb]{0.00,0.00,0.81}{#1}}
\newcommand{\FunctionTok}[1]{\textcolor[rgb]{0.13,0.29,0.53}{\textbf{#1}}}
\newcommand{\ImportTok}[1]{#1}
\newcommand{\InformationTok}[1]{\textcolor[rgb]{0.56,0.35,0.01}{\textbf{\textit{#1}}}}
\newcommand{\KeywordTok}[1]{\textcolor[rgb]{0.13,0.29,0.53}{\textbf{#1}}}
\newcommand{\NormalTok}[1]{#1}
\newcommand{\OperatorTok}[1]{\textcolor[rgb]{0.81,0.36,0.00}{\textbf{#1}}}
\newcommand{\OtherTok}[1]{\textcolor[rgb]{0.56,0.35,0.01}{#1}}
\newcommand{\PreprocessorTok}[1]{\textcolor[rgb]{0.56,0.35,0.01}{\textit{#1}}}
\newcommand{\RegionMarkerTok}[1]{#1}
\newcommand{\SpecialCharTok}[1]{\textcolor[rgb]{0.81,0.36,0.00}{\textbf{#1}}}
\newcommand{\SpecialStringTok}[1]{\textcolor[rgb]{0.31,0.60,0.02}{#1}}
\newcommand{\StringTok}[1]{\textcolor[rgb]{0.31,0.60,0.02}{#1}}
\newcommand{\VariableTok}[1]{\textcolor[rgb]{0.00,0.00,0.00}{#1}}
\newcommand{\VerbatimStringTok}[1]{\textcolor[rgb]{0.31,0.60,0.02}{#1}}
\newcommand{\WarningTok}[1]{\textcolor[rgb]{0.56,0.35,0.01}{\textbf{\textit{#1}}}}
\usepackage{longtable,booktabs,array}
\newcounter{none} % for unnumbered tables
\usepackage{calc} % for calculating minipage widths
% Correct order of tables after \paragraph or \subparagraph
\usepackage{etoolbox}
\makeatletter
\patchcmd\longtable{\par}{\if@noskipsec\mbox{}\fi\par}{}{}
\makeatother
% Allow footnotes in longtable head/foot
\IfFileExists{footnotehyper.sty}{\usepackage{footnotehyper}}{\usepackage{footnote}}
\makesavenoteenv{longtable}
\usepackage{graphicx}
\makeatletter
\newsavebox\pandoc@box
\newcommand*\pandocbounded[1]{% scales image to fit in text height/width
  \sbox\pandoc@box{#1}%
  \Gscale@div\@tempa{\textheight}{\dimexpr\ht\pandoc@box+\dp\pandoc@box\relax}%
  \Gscale@div\@tempb{\linewidth}{\wd\pandoc@box}%
  \ifdim\@tempb\p@<\@tempa\p@\let\@tempa\@tempb\fi% select the smaller of both
  \ifdim\@tempa\p@<\p@\scalebox{\@tempa}{\usebox\pandoc@box}%
  \else\usebox{\pandoc@box}%
  \fi%
}
% Set default figure placement to htbp
\def\fps@figure{htbp}
\makeatother
% definitions for citeproc citations
\NewDocumentCommand\citeproctext{}{}
\NewDocumentCommand\citeproc{mm}{%
  \begingroup\def\citeproctext{#2}\cite{#1}\endgroup}
\makeatletter
 % allow citations to break across lines
 \let\@cite@ofmt\@firstofone
 % avoid brackets around text for \cite:
 \def\@biblabel#1{}
 \def\@cite#1#2{{#1\if@tempswa , #2\fi}}
\makeatother
\newlength{\cslhangindent}
\setlength{\cslhangindent}{1.5em}
\newlength{\csllabelwidth}
\setlength{\csllabelwidth}{3em}
\newenvironment{CSLReferences}[2] % #1 hanging-indent, #2 entry-spacing
 {\begin{list}{}{%
  \setlength{\itemindent}{0pt}
  \setlength{\leftmargin}{0pt}
  \setlength{\parsep}{0pt}
  % turn on hanging indent if param 1 is 1
  \ifodd #1
   \setlength{\leftmargin}{\cslhangindent}
   \setlength{\itemindent}{-1\cslhangindent}
  \fi
  % set entry spacing
  \setlength{\itemsep}{#2\baselineskip}}}
 {\end{list}}
\usepackage{calc}
\newcommand{\CSLBlock}[1]{\hfill\break\parbox[t]{\linewidth}{\strut\ignorespaces#1\strut}}
\newcommand{\CSLLeftMargin}[1]{\parbox[t]{\csllabelwidth}{\strut#1\strut}}
\newcommand{\CSLRightInline}[1]{\parbox[t]{\linewidth - \csllabelwidth}{\strut#1\strut}}
\newcommand{\CSLIndent}[1]{\hspace{\cslhangindent}#1}
\ifLuaTeX
\usepackage[bidi=basic,shorthands=off]{babel}
\else
\usepackage[bidi=default,shorthands=off]{babel}
\fi
\ifLuaTeX
  \usepackage{selnolig} % disable illegal ligatures
\fi
\setlength{\emergencystretch}{3em} % prevent overfull lines
\providecommand{\tightlist}{%
  \setlength{\itemsep}{0pt}\setlength{\parskip}{0pt}}
\usepackage{bookmark}
\IfFileExists{xurl.sty}{\usepackage{xurl}}{} % add URL line breaks if available
\urlstyle{same}
% fallback for those not using the hyperref driver hyperxmp:
\makeatletter
\@ifundefined{xmpquote}{\newcommand{\xmpquote}[1]{#1}}{}
\makeatother
\hypersetup{
  pdftitle={Disease-modifying therapies and COVID-19 symptom reporting in multiple sclerosis},
  pdfauthor={Zina Tevzadze \textbar{} MSt Healthcare Data Science \textbar{} University of Cambridge},
  pdflang={en-GB},
  hidelinks,
  pdfcreator={LaTeX via pandoc}}

\title{Disease-modifying therapies and COVID-19 symptom reporting in
multiple sclerosis}
\usepackage{etoolbox}
\makeatletter
\providecommand{\subtitle}[1]{% add subtitle to \maketitle
  \apptocmd{\@title}{\par {\large #1 \par}}{}{}
}
\makeatother
\subtitle{Development and validation of an analysis-ready dataset}
\author{Zina Tevzadze \textbar{} MSt Healthcare Data Science \textbar{}
University of Cambridge}
\date{November 2026}

\begin{document}
\maketitle

{
\setcounter{tocdepth}{2}
\tableofcontents
}
\section*{Abstract}\label{abstract}
\addcontentsline{toc}{section}{Abstract}

\textbf{Background:} Disease-modifying therapies (DMTs) may alter immune
responses, but symptom reports and diagnostic categories are not
interchangeable (\citeproc{ref-sormani2021}{{Sormani et al.} 2021};
\citeproc{ref-struyf2022}{{Struyf et al.} 2022}). \textbf{Aim:} Develop
a documented, validated data asset for examining treatment--symptom
associations and reporting heterogeneity. \textbf{Methods:} Audit the
open GDSI dataset, specify phenotypes and assess readiness for logistic
regression. \textbf{Results:} To be completed after local execution and
interpretation. \textbf{Conclusion:} To be written independently after
validation.

\section{Introduction and Literature
Review}\label{introduction-and-literature-review}

\subsection{Clinical background}\label{clinical-background}

Multiple sclerosis (MS) is a chronic immune-mediated disorder of the
central nervous system characterised by inflammation, demyelination and
neurodegeneration. Disease-modifying therapies (DMTs) aim to reduce
inflammatory disease activity and limit disability accumulation.
However, their effects on immune function differ according to their
mechanisms of action, potentially influencing susceptibility to
infection and subsequent clinical manifestations.

The COVID-19 pandemic raised concerns regarding the management of people
with MS, particularly those receiving immunosuppressive or
immunomodulatory therapies. International data-sharing initiatives were
established to investigate these concerns using real-world clinical and
patient-reported information (\citeproc{ref-peeters2020}{{Peeters et
al.} 2020}).

\subsection{Disease-modifying therapies and
COVID-19}\label{disease-modifying-therapies-and-covid-19}

Observational studies have identified heterogeneous associations between
individual DMTs and COVID-19 severity. Sormani et al.~reported an
association between anti-CD20 treatment and more severe COVID-19
outcomes (\citeproc{ref-sormani2021}{{Sormani et al.} 2021}). Similarly,
international analyses identified associations between rituximab,
ocrelizumab and adverse clinical outcomes, alongside patient
characteristics such as age, MS phenotype and disability
(\citeproc{ref-simpson2022}{{Simpson-Yap et al.} 2022}).

These findings suggest that treating DMT exposure as a homogeneous
category may obscure clinically meaningful differences. In particular,
B-cell-depleting therapies have distinct immunological effects compared
with other DMT classes, motivating a mechanism-based classification of
treatment exposure.

Nevertheless, evidence concerning COVID-19 severity among infected
individuals cannot be directly extrapolated to the occurrence of
reported symptoms. Community-based research has demonstrated that
associations may differ according to whether COVID-19 is self-diagnosed
or laboratory-confirmed (\citeproc{ref-evangelou2021}{{Evangelou et al.}
2021}). Furthermore, symptoms alone have limited diagnostic specificity
for SARS-CoV-2 infection (\citeproc{ref-struyf2022}{{Struyf et al.}
2022}).

\subsection{Research rationale and
objectives}\label{research-rationale-and-objectives}

The publicly available COVID-MS dataset provides an opportunity to
investigate treatment--symptom associations using heterogeneous
real-world data (\citeproc{ref-khan2024}{Khan et al. 2024}). However,
differences in reporting source, diagnostic classification, clinical
granularity and data completeness may affect the validity and
interpretability of subsequent analyses.

\textbf{Research question:} How are disease-modifying therapy status and
anti-CD20 exposure associated with reported COVID-19-compatible symptoms
in people with multiple sclerosis?

This study aims to construct and validate an analysis-ready dataset
suitable for multivariable logistic regression. The approach emphasises
clinically justified exposure definitions, transparent missing-data
management, assessment of feature redundancy and evaluation of
statistical support.

Diagnostic and reporting heterogeneity will be examined to establish the
limitations of the resulting data asset and its suitability for the
intended analytical use.

\section{Methods}\label{methods}

\subsection{Data source and study
design}\label{data-source-and-study-design}

This study used the publicly available COVID-19 in MS Global Data
Sharing Initiative (GDSI) dataset, obtained through the Cambridge
Healthcare Data Science repository (\citeproc{ref-khan2024}{Khan et al.
2024}). Data were collected using a questionnaire-based platform through
which clinicians, people with multiple sclerosis (PwMS), or their
representatives submitted demographic, clinical, treatment and
COVID-19-related information.

The released dataset underwent de-identification, including age
categorisation, date perturbation, k-anonymity and l-diversity
procedures. These transformations preserve privacy but constrain
clinical granularity and temporal interpretation.

An observational, cross-sectional secondary analysis was specified, with
the individual source record as the unit of observation. The source
population was retained irrespective of COVID-19 diagnostic
classification, enabling investigation of reported symptoms without
restricting eligibility to suspected or confirmed infections.

Although source identifiers were unique, patient-level uniqueness across
reporting streams could not be independently verified.

\subsection{Variable selection and phenotype
definitions}\label{variable-selection-and-phenotype-definitions}

The primary outcome was defined as the reported presence of one or more
COVID-19-compatible symptoms, using the source variable
\texttt{covid19\_has\_symptoms}. Positive and negative responses were
encoded as 1 and 0, respectively, while missing responses remained
unavailable. This phenotype represents reported symptom occurrence
rather than confirmed SARS-CoV-2 infection or clinical severity.

DMT exposure was defined at two levels: treatment status at data entry
and mechanism-informed treatment groups. Individual drug identities were
preserved to maintain clinical granularity. Anti-CD20 therapies
(rituximab and ocrelizumab) were identified separately for comparison
with other identified current DMTs, informed by previous observational
evidence (\citeproc{ref-sormani2021}{{Sormani et al.} 2021};
\citeproc{ref-simpson2022}{{Simpson-Yap et al.} 2022}). Unspecified and
missing drug identities were distinguished from the absence of current
treatment and excluded from the anti-CD20 comparison.

Age, sex, MS phenotype and comorbidity status were selected as candidate
adjustment variables. Diagnostic classification and reporting source
were retained to evaluate reporting heterogeneity. Complete variable
definitions, derivation rules and analytical roles were documented in
the accompanying metadata.

Individual DMT identities were retained to preserve clinical
granularity. Mechanism-informed treatment groups were constructed for
descriptive analysis, with anti-CD20 exposure specified separately among
current users of identified therapies. Unspecified drug identities were
retained as a distinct category rather than assumed to represent
non-anti-CD20 treatment.

\begin{longtable}[]{@{}
  >{\raggedright\arraybackslash}p{(\linewidth - 4\tabcolsep) * \real{0.3826}}
  >{\raggedright\arraybackslash}p{(\linewidth - 4\tabcolsep) * \real{0.3652}}
  >{\raggedright\arraybackslash}p{(\linewidth - 4\tabcolsep) * \real{0.2522}}@{}}
\caption{Proposed analytical variable specification.}\tabularnewline
\toprule\noalign{}
\begin{minipage}[b]{\linewidth}\raggedright
Variable
\end{minipage} & \begin{minipage}[b]{\linewidth}\raggedright
Coding
\end{minipage} & \begin{minipage}[b]{\linewidth}\raggedright
Role
\end{minipage} \\
\midrule\noalign{}
\endfirsthead
\toprule\noalign{}
\begin{minipage}[b]{\linewidth}\raggedright
Variable
\end{minipage} & \begin{minipage}[b]{\linewidth}\raggedright
Coding
\end{minipage} & \begin{minipage}[b]{\linewidth}\raggedright
Role
\end{minipage} \\
\midrule\noalign{}
\endhead
\bottomrule\noalign{}
\endlastfoot
DMT status & Never / Not currently / Current & Primary exposure \\
DMT class & Mechanism-informed treatment groups & Descriptive
classification \\
Anti-CD20 exposure & Anti-CD20 / Other identified DMT & Secondary
comparison \\
COVID-19-compatible symptoms & 0 = No; 1 = Yes & Primary outcome \\
Age & Categorical age at reporting & Candidate adjustment \\
Sex & Source-recorded categories & Candidate adjustment \\
MS phenotype & Relapsing-remitting / Progressive / Other & Candidate
adjustment \\
Comorbidity & Yes / No & Candidate adjustment \\
Reporting source & Patient / Clinician & Reporting heterogeneity \\
COVID-19 diagnosis & Not suspected / Suspected / Confirmed & Diagnostic
heterogeneity \\
Hospitalisation, ICU admission, ventilation & Source-recorded outcomes &
Supporting clinical outcomes \\
\end{longtable}

\subsection{Data quality assessment and
preprocessing}\label{data-quality-assessment-and-preprocessing}

Source-data quality was assessed across completeness, validity,
uniqueness and cross-field consistency. Structural checks included
record identifiers, dataset dimensions, permitted categorical values and
agreement between related clinical variables.

Missingness was quantified overall and by reporting source. Conditional
relationships between variables were examined to distinguish
structurally inapplicable information from unavailable clinical
information. Missing values were not automatically interpreted as
negative findings or subjected to blanket imputation.

Categorical variables were harmonised into explicitly defined levels,
preserving the clinical meaning of the original source categories.
Treatment identities were retained alongside derived treatment groups.

All source records were retained in the master ARD. Analysis-specific
eligibility indicators were derived according to outcome availability,
exposure definition and covariate completeness. Variables with
insufficient observed variation, uncertain clinical interpretation or
inadequate completeness were excluded from the primary analytical
specification.

The resulting dataset was subjected to automated structural and logical
validation, including verification of derived treatment phenotypes,
eligibility criteria and consistency between the in-memory dataset and
its exported RDS representation.

\begin{longtable}[]{@{}
  >{\raggedright\arraybackslash}p{(\linewidth - 2\tabcolsep) * \real{0.3117}}
  >{\raggedright\arraybackslash}p{(\linewidth - 2\tabcolsep) * \real{0.6883}}@{}}
\caption{Summary of data preparation procedures.}\tabularnewline
\toprule\noalign{}
\begin{minipage}[b]{\linewidth}\raggedright
Operation
\end{minipage} & \begin{minipage}[b]{\linewidth}\raggedright
Implementation
\end{minipage} \\
\midrule\noalign{}
\endfirsthead
\toprule\noalign{}
\begin{minipage}[b]{\linewidth}\raggedright
Operation
\end{minipage} & \begin{minipage}[b]{\linewidth}\raggedright
Implementation
\end{minipage} \\
\midrule\noalign{}
\endhead
\bottomrule\noalign{}
\endlastfoot
Structural validation & Dimensions, identifiers and permitted values \\
Semantic harmonisation & Explicit clinical categories and reference
levels \\
Phenotype engineering & Binary symptoms, DMT groups and anti-CD20
exposure \\
Missing-data management & Preservation of missing and inapplicable
information \\
Cohort construction & Analysis-specific eligibility indicators \\
Export validation & Consistency between in-memory and exported RDS
data \\
\end{longtable}

\subsection{Analytical specification and
validation}\label{analytical-specification-and-validation}

Multivariable binary logistic regression was selected as the intended
analytical method, reflecting the binary nature of the symptom outcome.
Two exposure specifications were defined: overall DMT status and
anti-CD20 treatment compared with other identified current therapies.

Age, sex, MS phenotype and comorbidity status were specified as
candidate adjustment variables based on clinical relevance. Treatment
status and treatment class were not included simultaneously because of
their definitional dependence.

Logistic regression permits estimation of adjusted odds ratios for
categorical treatment exposures while accommodating demographic and
clinical covariates. Given the explanatory objective, emphasis was
placed on exposure definition, coefficient estimability and potential
confounding rather than predictive performance.

Analytical readiness was assessed through cohort reconciliation, outcome
variation, exposure-specific event counts and model-matrix rank. Further
assessment of multicollinearity and separation was specified to evaluate
the feasibility of conventional logistic regression.

Diagnostic classification and reporting source were examined
descriptively to characterise reporting heterogeneity and assess the
feasibility of subgroup analyses. The intended analyses were
associational rather than causal, reflecting the cross-sectional study
design and limitations of the source data.

Exposure-specific symptom proportions and 95\% Wilson confidence
intervals were calculated descriptively to assess outcome variation and
statistical support across treatment groups.

The general analytical specification was:

\[
\log\left(\frac{p_i}{1-p_i}\right)
=
\beta_0 +
\boldsymbol{\beta}^{T}\mathbf{D}_i +
\boldsymbol{\gamma}^{T}\mathbf{X}_i
\]

where \(p_i\) denotes the probability of reported COVID-19-compatible
symptoms, \(\mathbf{D}_i\) represents the treatment exposure and
\(\mathbf{X}_i\) contains the candidate adjustment variables.

\section{Results}\label{results}

\subsection{Source-data quality and
missingness}\label{source-data-quality-and-missingness}

The source dataset comprised 1141 records and 47 variables. No missing
or duplicated source record identifiers were identified.

Patient-reported records accounted for 92.6\% of the dataset, compared
with 7.4\% submitted by clinicians.

\begin{Shaded}
\begin{Highlighting}[]
\NormalTok{missing\_labels }\OtherTok{\textless{}{-}} \FunctionTok{c}\NormalTok{(}
  \AttributeTok{edss\_in\_cat2 =} \StringTok{"EDSS"}\NormalTok{,}
  \AttributeTok{dmt\_type\_overall =} \StringTok{"DMT identity"}\NormalTok{,}
  \AttributeTok{current\_or\_former\_smoker =} \StringTok{"Smoking history"}\NormalTok{,}
  \AttributeTok{covid19\_has\_symptoms =} \StringTok{"Symptom outcome"}\NormalTok{,}
  \AttributeTok{bmi\_in\_cat2 =} \StringTok{"BMI"}
\NormalTok{)}

\NormalTok{missing\_by\_source }\SpecialCharTok{|\textgreater{}}
  \FunctionTok{filter}\NormalTok{(variable }\SpecialCharTok{\%in\%} \FunctionTok{names}\NormalTok{(missing\_labels)) }\SpecialCharTok{|\textgreater{}}
  \FunctionTok{mutate}\NormalTok{(}
    \AttributeTok{variable\_label =} \FunctionTok{factor}\NormalTok{(}
      \FunctionTok{unname}\NormalTok{(missing\_labels[variable]),}
      \AttributeTok{levels =} \FunctionTok{rev}\NormalTok{(}\FunctionTok{unname}\NormalTok{(missing\_labels))}
\NormalTok{    ),}
    \AttributeTok{report\_source =} \FunctionTok{factor}\NormalTok{(}
\NormalTok{      report\_source,}
      \AttributeTok{levels =} \FunctionTok{c}\NormalTok{(}\StringTok{"clinicians"}\NormalTok{, }\StringTok{"patients"}\NormalTok{),}
      \AttributeTok{labels =} \FunctionTok{c}\NormalTok{(}
        \StringTok{"Clinician (n = 84)"}\NormalTok{,}
        \StringTok{"Patient (n = 1,057)"}
\NormalTok{      )}
\NormalTok{    )}
\NormalTok{  ) }\SpecialCharTok{|\textgreater{}}
  \FunctionTok{ggplot}\NormalTok{(}
    \FunctionTok{aes}\NormalTok{(}
      \AttributeTok{x =}\NormalTok{ report\_source,}
      \AttributeTok{y =}\NormalTok{ variable\_label,}
      \AttributeTok{fill =}\NormalTok{ missing\_pct}
\NormalTok{    )}
\NormalTok{  ) }\SpecialCharTok{+}
  \FunctionTok{geom\_tile}\NormalTok{(}
    \AttributeTok{colour =} \StringTok{"white"}\NormalTok{,}
    \AttributeTok{linewidth =} \DecValTok{1}
\NormalTok{  ) }\SpecialCharTok{+}
  \FunctionTok{geom\_text}\NormalTok{(}
    \FunctionTok{aes}\NormalTok{(}
      \AttributeTok{label =} \FunctionTok{sprintf}\NormalTok{(}\StringTok{"\%.1f\%\%"}\NormalTok{, missing\_pct),}
      \AttributeTok{colour =}\NormalTok{ missing\_pct }\SpecialCharTok{\textgreater{}=} \DecValTok{40}
\NormalTok{    ),}
    \AttributeTok{size =} \DecValTok{4}
\NormalTok{  ) }\SpecialCharTok{+}
  \FunctionTok{scale\_fill\_gradientn}\NormalTok{(}
    \AttributeTok{colours =} \FunctionTok{rev}\NormalTok{(rocket),}
    \AttributeTok{limits =} \FunctionTok{c}\NormalTok{(}\DecValTok{0}\NormalTok{, }\DecValTok{100}\NormalTok{),}
    \AttributeTok{breaks =} \FunctionTok{seq}\NormalTok{(}\DecValTok{0}\NormalTok{, }\DecValTok{100}\NormalTok{, }\DecValTok{25}\NormalTok{)}
\NormalTok{  ) }\SpecialCharTok{+}
  \FunctionTok{scale\_colour\_manual}\NormalTok{(}
    \AttributeTok{values =} \FunctionTok{c}\NormalTok{(}
      \StringTok{"TRUE"} \OtherTok{=} \StringTok{"white"}\NormalTok{,}
      \StringTok{"FALSE"} \OtherTok{=} \StringTok{"\#35193E"}
\NormalTok{    ),}
    \AttributeTok{guide =} \StringTok{"none"}
\NormalTok{  ) }\SpecialCharTok{+}
  \FunctionTok{labs}\NormalTok{(}
    \AttributeTok{title =} \StringTok{"Missingness by reporting source"}\NormalTok{,}
    \AttributeTok{subtitle =} \StringTok{"Selected clinically relevant variables"}\NormalTok{,}
    \AttributeTok{x =} \ConstantTok{NULL}\NormalTok{,}
    \AttributeTok{y =} \ConstantTok{NULL}\NormalTok{,}
    \AttributeTok{fill =} \StringTok{"Missing (\%)"}
\NormalTok{  ) }\SpecialCharTok{+}
  \FunctionTok{theme}\NormalTok{(}
    \AttributeTok{panel.grid =} \FunctionTok{element\_blank}\NormalTok{(),}
    \AttributeTok{axis.ticks =} \FunctionTok{element\_blank}\NormalTok{(),}
    \AttributeTok{legend.position =} \StringTok{"bottom"}
\NormalTok{  )}
\end{Highlighting}
\end{Shaded}

\begin{figure}

{\centering \includegraphics[width=1\linewidth]{report_files/figure-latex/missing-source-1} 

}

\caption{Missingness in selected clinical variables by reporting source.}\label{fig:missing-source}
\end{figure}

Missingness varied substantially by reporting source. EDSS was more
frequently unavailable in patient reports than clinician reports (57.4\%
versus 13.1\%), whereas BMI and smoking history exhibited greater
missingness in clinician reports. All eight missing primary symptom
outcomes occurred among clinician-reported records. These differences
demonstrate heterogeneous data availability across reporting streams and
highlight the potential for selection bias under complete-case analysis.

\begin{longtable}[]{@{}lrr@{}}
\caption{Availability of drug identity by treatment
status.}\tabularnewline
\toprule\noalign{}
Treatment status & Records & Missing DMT identity \\
\midrule\noalign{}
\endfirsthead
\toprule\noalign{}
Treatment status & Records & Missing DMT identity \\
\midrule\noalign{}
\endhead
\bottomrule\noalign{}
\endlastfoot
never\_treated & 160 & 160 \\
no & 101 & 0 \\
yes & 880 & 1 \\
\end{longtable}

\begin{longtable}[]{@{}
  >{\raggedright\arraybackslash}p{(\linewidth - 4\tabcolsep) * \real{0.2468}}
  >{\raggedleft\arraybackslash}p{(\linewidth - 4\tabcolsep) * \real{0.1558}}
  >{\raggedright\arraybackslash}p{(\linewidth - 4\tabcolsep) * \real{0.5974}}@{}}
\caption{Selected exclusions from the primary analytical
specification.}\tabularnewline
\toprule\noalign{}
\begin{minipage}[b]{\linewidth}\raggedright
Variable
\end{minipage} & \begin{minipage}[b]{\linewidth}\raggedleft
Missing (\%)
\end{minipage} & \begin{minipage}[b]{\linewidth}\raggedright
Exclusion rationale
\end{minipage} \\
\midrule\noalign{}
\endfirsthead
\toprule\noalign{}
\begin{minipage}[b]{\linewidth}\raggedright
Variable
\end{minipage} & \begin{minipage}[b]{\linewidth}\raggedleft
Missing (\%)
\end{minipage} & \begin{minipage}[b]{\linewidth}\raggedright
Exclusion rationale
\end{minipage} \\
\midrule\noalign{}
\endhead
\bottomrule\noalign{}
\endlastfoot
Treatment duration & 80.5 & Limited completeness and temporal
information \\
Smoking history & 56.6 & Incomplete and ambiguous coding \\
EDSS & 54.2 & No observed variation \\
BMI & 22.0 & Incomplete and coarsely categorised \\
\end{longtable}

Variable exclusion was not determined by missingness alone. All 523
available EDSS observations belonged to the same released disability
category, precluding estimation of associations with disability level.
Exclusion of potentially relevant covariates, particularly disability
and BMI, nevertheless limits adjustment for clinical differences between
treatment groups.

\subsection{Feature engineering and cohort
construction}\label{feature-engineering-and-cohort-construction}

The preparation pipeline produced an ARD containing 1141 records and 20
variables. All source records were retained.

The resulting data structure incorporated harmonised clinical variables,
a binary symptom phenotype, derived treatment classifications and
analysis-specific eligibility indicators. Original drug identities were
preserved alongside the derived exposure definitions.

\begin{longtable}[]{@{}
  >{\raggedright\arraybackslash}p{(\linewidth - 6\tabcolsep) * \real{0.3944}}
  >{\raggedleft\arraybackslash}p{(\linewidth - 6\tabcolsep) * \real{0.1127}}
  >{\raggedright\arraybackslash}p{(\linewidth - 6\tabcolsep) * \real{0.3380}}
  >{\raggedleft\arraybackslash}p{(\linewidth - 6\tabcolsep) * \real{0.1549}}@{}}
\caption{Derived treatment classification among current DMT
users.}\tabularnewline
\toprule\noalign{}
\begin{minipage}[b]{\linewidth}\raggedright
Treatment group
\end{minipage} & \begin{minipage}[b]{\linewidth}\raggedleft
Records
\end{minipage} & \begin{minipage}[b]{\linewidth}\raggedright
Treatments
\end{minipage} & \begin{minipage}[b]{\linewidth}\raggedleft
Percentage
\end{minipage} \\
\midrule\noalign{}
\endfirsthead
\toprule\noalign{}
\begin{minipage}[b]{\linewidth}\raggedright
Treatment group
\end{minipage} & \begin{minipage}[b]{\linewidth}\raggedleft
Records
\end{minipage} & \begin{minipage}[b]{\linewidth}\raggedright
Treatments
\end{minipage} & \begin{minipage}[b]{\linewidth}\raggedleft
Percentage
\end{minipage} \\
\midrule\noalign{}
\endhead
\bottomrule\noalign{}
\endlastfoot
Other unspecified & 185 & another drug not listed & 21.0 \\
Injectable immunomodulators & 158 & glatiramer, interferon & 18.0 \\
Fumarates & 145 & dimethyl fumarate & 16.5 \\
S1P modulators & 116 & fingolimod & 13.2 \\
Anti-CD20 & 99 & ocrelizumab, rituximab & 11.2 \\
Trafficking inhibition & 65 & natalizumab & 7.4 \\
Teriflunomide & 63 & teriflunomide & 7.2 \\
Immune reconstitution & 48 & alemtuzumab, cladribine & 5.5 \\
Unknown drug & 1 & Identity unavailable & 0.1 \\
\end{longtable}

\begin{Shaded}
\begin{Highlighting}[]
\NormalTok{ard }\SpecialCharTok{|\textgreater{}}
  \FunctionTok{filter}\NormalTok{(dmt\_status }\SpecialCharTok{==} \StringTok{"Current"}\NormalTok{) }\SpecialCharTok{|\textgreater{}}
  \FunctionTok{count}\NormalTok{(dmt\_class, }\AttributeTok{name =} \StringTok{"records"}\NormalTok{) }\SpecialCharTok{|\textgreater{}}
  \FunctionTok{mutate}\NormalTok{(}
    \AttributeTok{group =} \FunctionTok{case\_when}\NormalTok{(}
\NormalTok{      dmt\_class }\SpecialCharTok{==} \StringTok{"Anti{-}CD20"} \SpecialCharTok{\textasciitilde{}} \StringTok{"Target exposure"}\NormalTok{,}
\NormalTok{      dmt\_class }\SpecialCharTok{\%in\%} \FunctionTok{c}\NormalTok{(}
        \StringTok{"Other unspecified"}\NormalTok{, }\StringTok{"Unknown drug"}
\NormalTok{      ) }\SpecialCharTok{\textasciitilde{}} \StringTok{"Unidentified treatment"}\NormalTok{,}
      \ConstantTok{TRUE} \SpecialCharTok{\textasciitilde{}} \StringTok{"Other identified treatment"}
\NormalTok{    )}
\NormalTok{  ) }\SpecialCharTok{|\textgreater{}}
  \FunctionTok{ggplot}\NormalTok{(}
    \FunctionTok{aes}\NormalTok{(}
      \AttributeTok{x =} \FunctionTok{reorder}\NormalTok{(dmt\_class, records),}
      \AttributeTok{y =}\NormalTok{ records,}
      \AttributeTok{fill =}\NormalTok{ group}
\NormalTok{    )}
\NormalTok{  ) }\SpecialCharTok{+}
  \FunctionTok{geom\_col}\NormalTok{(}\AttributeTok{width =} \FloatTok{0.7}\NormalTok{) }\SpecialCharTok{+}
  \FunctionTok{geom\_text}\NormalTok{(}
    \FunctionTok{aes}\NormalTok{(}\AttributeTok{label =}\NormalTok{ records),}
    \AttributeTok{hjust =} \SpecialCharTok{{-}}\FloatTok{0.2}\NormalTok{,}
    \AttributeTok{size =} \FloatTok{3.5}
\NormalTok{  ) }\SpecialCharTok{+}
  \FunctionTok{coord\_flip}\NormalTok{() }\SpecialCharTok{+}
  \FunctionTok{scale\_fill\_manual}\NormalTok{(}
    \AttributeTok{values =} \FunctionTok{c}\NormalTok{(}
      \StringTok{"Target exposure"} \OtherTok{=} \StringTok{"\#AD1759"}\NormalTok{,}
      \StringTok{"Other identified treatment"} \OtherTok{=} \StringTok{"\#35193E"}\NormalTok{,}
      \StringTok{"Unidentified treatment"} \OtherTok{=} \StringTok{"\#B9B2BB"}
\NormalTok{    )}
\NormalTok{  ) }\SpecialCharTok{+}
  \FunctionTok{scale\_y\_continuous}\NormalTok{(}
    \AttributeTok{expand =} \FunctionTok{expansion}\NormalTok{(}\AttributeTok{mult =} \FunctionTok{c}\NormalTok{(}\DecValTok{0}\NormalTok{, }\FloatTok{0.15}\NormalTok{))}
\NormalTok{  ) }\SpecialCharTok{+}
  \FunctionTok{labs}\NormalTok{(}
    \AttributeTok{title =} \StringTok{"Distribution of current DMT users"}\NormalTok{,}
    \AttributeTok{subtitle =} \StringTok{"Mechanism{-}informed treatment classification"}\NormalTok{,}
    \AttributeTok{x =} \ConstantTok{NULL}\NormalTok{,}
    \AttributeTok{y =} \StringTok{"Number of records"}\NormalTok{,}
    \AttributeTok{fill =} \ConstantTok{NULL}
\NormalTok{  )}
\end{Highlighting}
\end{Shaded}

\begin{figure}

{\centering \includegraphics[width=1\linewidth]{report_files/figure-latex/dmt-distribution-1} 

}

\caption{Distribution of derived treatment groups among current DMT users.}\label{fig:dmt-distribution}
\end{figure}

\begin{longtable}[]{@{}lrr@{}}
\caption{Record retention and analytical eligibility.}\tabularnewline
\toprule\noalign{}
Analytical cohort & Records & Source records (\%) \\
\midrule\noalign{}
\endfirsthead
\toprule\noalign{}
Analytical cohort & Records & Source records (\%) \\
\midrule\noalign{}
\endhead
\bottomrule\noalign{}
\endlastfoot
Source records & 1141 & 100.0 \\
Primary eligible & 1133 & 99.3 \\
Adjusted eligible & 1133 & 99.3 \\
Diagnostic subset & 275 & 24.1 \\
Secondary anti-CD20 cohort & 687 & 60.2 \\
\end{longtable}

The primary adjusted cohort retained 99.3\% of source records, with
eight exclusions attributable to missing symptom outcomes. Inclusion of
the four candidate adjustment variables resulted in no additional record
loss. Among current users of identified therapies, 687 records satisfied
secondary analytical eligibility, including 99 anti-CD20 and 588
comparator records.

\subsection{Analytical readiness and exploratory
findings}\label{analytical-readiness-and-exploratory-findings}

\subsubsection{Diagnostic and reporting
heterogeneity}\label{diagnostic-and-reporting-heterogeneity}

The relationship between diagnostic classification and reported symptoms
was examined to assess outcome interpretability and the feasibility of
restricting the analytical population to suspected or confirmed COVID-19
cases.

\begin{Shaded}
\begin{Highlighting}[]
\NormalTok{diagnosis\_summary }\OtherTok{\textless{}{-}}\NormalTok{ ard }\SpecialCharTok{|\textgreater{}}
  \FunctionTok{mutate}\NormalTok{(}
    \AttributeTok{symptom\_status =} \FunctionTok{case\_when}\NormalTok{(}
\NormalTok{      symptoms }\SpecialCharTok{==} \DecValTok{1} \SpecialCharTok{\textasciitilde{}} \StringTok{"Symptoms"}\NormalTok{,}
\NormalTok{      symptoms }\SpecialCharTok{==} \DecValTok{0} \SpecialCharTok{\textasciitilde{}} \StringTok{"No symptoms"}\NormalTok{,}
      \ConstantTok{TRUE} \SpecialCharTok{\textasciitilde{}} \StringTok{"Missing"}
\NormalTok{    ),}
    \AttributeTok{symptom\_status =} \FunctionTok{factor}\NormalTok{(}
\NormalTok{      symptom\_status,}
      \AttributeTok{levels =} \FunctionTok{c}\NormalTok{(}
        \StringTok{"Symptoms"}\NormalTok{,}
        \StringTok{"No symptoms"}\NormalTok{,}
        \StringTok{"Missing"}
\NormalTok{      )}
\NormalTok{    )}
\NormalTok{  ) }\SpecialCharTok{|\textgreater{}}
  \FunctionTok{count}\NormalTok{(diagnosis, symptom\_status) }\SpecialCharTok{|\textgreater{}}
  \FunctionTok{group\_by}\NormalTok{(diagnosis) }\SpecialCharTok{|\textgreater{}}
  \FunctionTok{mutate}\NormalTok{(}
    \AttributeTok{proportion =}\NormalTok{ n }\SpecialCharTok{/} \FunctionTok{sum}\NormalTok{(n)}
\NormalTok{  ) }\SpecialCharTok{|\textgreater{}}
  \FunctionTok{ungroup}\NormalTok{()}

\NormalTok{diagnosis\_summary }\SpecialCharTok{|\textgreater{}}
  \FunctionTok{ggplot}\NormalTok{(}
    \FunctionTok{aes}\NormalTok{(}
      \AttributeTok{x =}\NormalTok{ diagnosis,}
      \AttributeTok{y =}\NormalTok{ proportion,}
      \AttributeTok{fill =}\NormalTok{ symptom\_status}
\NormalTok{    )}
\NormalTok{  ) }\SpecialCharTok{+}
  \FunctionTok{geom\_col}\NormalTok{(}
    \AttributeTok{width =} \FloatTok{0.65}\NormalTok{,}
    \AttributeTok{position =} \FunctionTok{position\_stack}\NormalTok{(}\AttributeTok{reverse =} \ConstantTok{TRUE}\NormalTok{)}
\NormalTok{  ) }\SpecialCharTok{+}
  \FunctionTok{geom\_text}\NormalTok{(}
    \FunctionTok{aes}\NormalTok{(}
      \AttributeTok{label =} \FunctionTok{if\_else}\NormalTok{(}
\NormalTok{        proportion }\SpecialCharTok{\textgreater{}=} \FloatTok{0.15}\NormalTok{,}
\NormalTok{        scales}\SpecialCharTok{::}\FunctionTok{percent}\NormalTok{(proportion, }\AttributeTok{accuracy =} \FloatTok{0.1}\NormalTok{),}
        \StringTok{""}
\NormalTok{      )}
\NormalTok{    ),}
    \AttributeTok{position =} \FunctionTok{position\_stack}\NormalTok{(}
      \AttributeTok{vjust =} \FloatTok{0.5}\NormalTok{,}
      \AttributeTok{reverse =} \ConstantTok{TRUE}
\NormalTok{    ),}
    \AttributeTok{colour =} \StringTok{"white"}\NormalTok{,}
    \AttributeTok{size =} \FloatTok{3.5}
\NormalTok{  ) }\SpecialCharTok{+}
  \FunctionTok{coord\_flip}\NormalTok{() }\SpecialCharTok{+}
  \FunctionTok{scale\_fill\_manual}\NormalTok{(}
    \AttributeTok{values =} \FunctionTok{c}\NormalTok{(}
      \StringTok{"Symptoms"} \OtherTok{=} \StringTok{"\#AD1759"}\NormalTok{,}
      \StringTok{"No symptoms"} \OtherTok{=} \StringTok{"\#35193E"}\NormalTok{,}
      \StringTok{"Missing"} \OtherTok{=} \StringTok{"\#D9D9D9"}
\NormalTok{    )}
\NormalTok{  ) }\SpecialCharTok{+}
  \FunctionTok{scale\_y\_continuous}\NormalTok{(}
    \AttributeTok{labels =}\NormalTok{ scales}\SpecialCharTok{::}\FunctionTok{label\_percent}\NormalTok{()}
\NormalTok{  ) }\SpecialCharTok{+}
  \FunctionTok{scale\_x\_discrete}\NormalTok{(}
    \AttributeTok{labels =} \FunctionTok{c}\NormalTok{(}
      \StringTok{"confirmed"} \OtherTok{=} \StringTok{"Confirmed"}\NormalTok{,}
      \StringTok{"suspected"} \OtherTok{=} \StringTok{"Suspected"}\NormalTok{,}
      \StringTok{"not\_suspected"} \OtherTok{=} \StringTok{"Not suspected"}
\NormalTok{    )}
\NormalTok{  ) }\SpecialCharTok{+}
  \FunctionTok{labs}\NormalTok{(}
    \AttributeTok{title =} \StringTok{"Symptom reporting by diagnosis"}\NormalTok{,}
    \AttributeTok{subtitle =} \StringTok{"Within{-}category outcome distributions"}\NormalTok{,}
    \AttributeTok{x =} \ConstantTok{NULL}\NormalTok{,}
    \AttributeTok{y =} \StringTok{"Proportion of records"}\NormalTok{,}
    \AttributeTok{fill =} \ConstantTok{NULL}
\NormalTok{  )}
\end{Highlighting}
\end{Shaded}

\begin{figure}

{\centering \includegraphics[width=1\linewidth]{report_files/figure-latex/diagnosis-symptoms-1} 

}

\caption{Distribution of reported symptoms within COVID-19 diagnostic categories.}\label{fig:diagnosis-symptoms}
\end{figure}

Symptom reporting differed markedly across diagnostic categories. All
suspected cases with observed symptom outcomes were symptomatic, while
89 records classified as not suspected also reported symptoms.
Restriction to suspected and confirmed cases would retain only four
asymptomatic records, substantially limiting outcome variation. These
findings support retaining all diagnostic categories in the primary
analytical cohort and treating diagnosis as a descriptive characteristic
rather than a routine adjustment variable.

\begin{Shaded}
\begin{Highlighting}[]
\NormalTok{ard }\SpecialCharTok{|\textgreater{}}
  \FunctionTok{count}\NormalTok{(reporting\_source, diagnosis) }\SpecialCharTok{|\textgreater{}}
  \FunctionTok{group\_by}\NormalTok{(reporting\_source) }\SpecialCharTok{|\textgreater{}}
  \FunctionTok{mutate}\NormalTok{(}
    \AttributeTok{proportion =}\NormalTok{ n }\SpecialCharTok{/} \FunctionTok{sum}\NormalTok{(n)}
\NormalTok{  ) }\SpecialCharTok{|\textgreater{}}
  \FunctionTok{ungroup}\NormalTok{() }\SpecialCharTok{|\textgreater{}}
  \FunctionTok{ggplot}\NormalTok{(}
    \FunctionTok{aes}\NormalTok{(}
      \AttributeTok{x =}\NormalTok{ reporting\_source,}
      \AttributeTok{y =}\NormalTok{ proportion,}
      \AttributeTok{fill =}\NormalTok{ diagnosis}
\NormalTok{    )}
\NormalTok{  ) }\SpecialCharTok{+}
  \FunctionTok{geom\_col}\NormalTok{(}\AttributeTok{width =} \FloatTok{0.6}\NormalTok{) }\SpecialCharTok{+}
  \FunctionTok{geom\_text}\NormalTok{(}
    \FunctionTok{aes}\NormalTok{(}
      \AttributeTok{label =} \FunctionTok{if\_else}\NormalTok{(}
\NormalTok{        proportion }\SpecialCharTok{\textgreater{}=} \FloatTok{0.12}\NormalTok{,}
\NormalTok{        scales}\SpecialCharTok{::}\FunctionTok{percent}\NormalTok{(proportion, }\AttributeTok{accuracy =} \FloatTok{0.1}\NormalTok{),}
        \StringTok{""}
\NormalTok{      ),}
      \AttributeTok{colour =}\NormalTok{ diagnosis }\SpecialCharTok{==} \StringTok{"suspected"}
\NormalTok{    ),}
    \AttributeTok{position =} \FunctionTok{position\_stack}\NormalTok{(}\AttributeTok{vjust =} \FloatTok{0.5}\NormalTok{),}
    \AttributeTok{size =} \FloatTok{3.5}
\NormalTok{  ) }\SpecialCharTok{+}
  \FunctionTok{coord\_flip}\NormalTok{() }\SpecialCharTok{+}
  \FunctionTok{scale\_fill\_manual}\NormalTok{(}
    \AttributeTok{values =} \FunctionTok{c}\NormalTok{(}
      \StringTok{"confirmed"} \OtherTok{=} \StringTok{"\#AD1759"}\NormalTok{,}
      \StringTok{"suspected"} \OtherTok{=} \StringTok{"\#F58860"}\NormalTok{,}
      \StringTok{"not\_suspected"} \OtherTok{=} \StringTok{"\#35193E"}
\NormalTok{    ),}
    \AttributeTok{labels =} \FunctionTok{c}\NormalTok{(}
      \StringTok{"confirmed"} \OtherTok{=} \StringTok{"Confirmed"}\NormalTok{,}
      \StringTok{"suspected"} \OtherTok{=} \StringTok{"Suspected"}\NormalTok{,}
      \StringTok{"not\_suspected"} \OtherTok{=} \StringTok{"Not suspected"}
\NormalTok{    )}
\NormalTok{  ) }\SpecialCharTok{+}
  \FunctionTok{scale\_colour\_manual}\NormalTok{(}
    \AttributeTok{values =} \FunctionTok{c}\NormalTok{(}
      \StringTok{"TRUE"} \OtherTok{=} \StringTok{"\#35193E"}\NormalTok{,}
      \StringTok{"FALSE"} \OtherTok{=} \StringTok{"white"}
\NormalTok{    ),}
    \AttributeTok{guide =} \StringTok{"none"}
\NormalTok{  ) }\SpecialCharTok{+}
  \FunctionTok{scale\_y\_continuous}\NormalTok{(}
    \AttributeTok{labels =}\NormalTok{ scales}\SpecialCharTok{::}\FunctionTok{label\_percent}\NormalTok{()}
\NormalTok{  ) }\SpecialCharTok{+}
  \FunctionTok{scale\_x\_discrete}\NormalTok{(}
    \AttributeTok{labels =} \FunctionTok{c}\NormalTok{(}
      \StringTok{"clinicians"} \OtherTok{=} \StringTok{"Clinicians"}\NormalTok{,}
      \StringTok{"patients"} \OtherTok{=} \StringTok{"Patients"}
\NormalTok{    )}
\NormalTok{  ) }\SpecialCharTok{+}
  \FunctionTok{labs}\NormalTok{(}
    \AttributeTok{title =} \StringTok{"Diagnosis by reporting source"}\NormalTok{,}
    \AttributeTok{subtitle =} \StringTok{"Within{-}source diagnostic distributions"}\NormalTok{,}
    \AttributeTok{x =} \ConstantTok{NULL}\NormalTok{,}
    \AttributeTok{y =} \StringTok{"Proportion of records"}\NormalTok{,}
    \AttributeTok{fill =} \ConstantTok{NULL}
\NormalTok{  )}
\end{Highlighting}
\end{Shaded}

\begin{figure}

{\centering \includegraphics[width=1\linewidth]{report_files/figure-latex/reporting-diagnosis-1} 

}

\caption{Diagnostic composition of clinician- and patient-reported records.}\label{fig:reporting-diagnosis}
\end{figure}

Confirmed COVID-19 accounted for 52.4\% of clinician reports but only
1.5\% of patient reports. This substantial difference suggests
heterogeneous case ascertainment and reporting processes. Consequently,
reporting source may capture differences in selection, diagnostic
verification and clinical presentation, rather than functioning solely
as a conventional demographic covariate.

\subsubsection{Exposure-specific statistical
support}\label{exposure-specific-statistical-support}

\begin{Shaded}
\begin{Highlighting}[]
\NormalTok{primary\_support }\OtherTok{\textless{}{-}}\NormalTok{ ard }\SpecialCharTok{|\textgreater{}}
  \FunctionTok{filter}\NormalTok{(eligible\_adjusted) }\SpecialCharTok{|\textgreater{}}
  \FunctionTok{group\_by}\NormalTok{(dmt\_status) }\SpecialCharTok{|\textgreater{}}
  \FunctionTok{summarise}\NormalTok{(}
    \AttributeTok{records =} \FunctionTok{n}\NormalTok{(),}
    \AttributeTok{events =} \FunctionTok{sum}\NormalTok{(symptoms }\SpecialCharTok{==} \DecValTok{1}\NormalTok{),}
    \AttributeTok{non\_events =} \FunctionTok{sum}\NormalTok{(symptoms }\SpecialCharTok{==} \DecValTok{0}\NormalTok{),}
    \AttributeTok{.groups =} \StringTok{"drop"}
\NormalTok{  ) }\SpecialCharTok{|\textgreater{}}
  \FunctionTok{mutate}\NormalTok{(}
    \AttributeTok{proportion =}\NormalTok{ events }\SpecialCharTok{/}\NormalTok{ records,}
    \AttributeTok{z =} \FunctionTok{qnorm}\NormalTok{(}\FloatTok{0.975}\NormalTok{),}

    \AttributeTok{denominator =} \DecValTok{1} \SpecialCharTok{+}\NormalTok{ z}\SpecialCharTok{\^{}}\DecValTok{2} \SpecialCharTok{/}\NormalTok{ records,}

    \AttributeTok{centre =}\NormalTok{ (}
\NormalTok{      proportion }\SpecialCharTok{+}\NormalTok{ z}\SpecialCharTok{\^{}}\DecValTok{2} \SpecialCharTok{/}\NormalTok{ (}\DecValTok{2} \SpecialCharTok{*}\NormalTok{ records)}
\NormalTok{    ) }\SpecialCharTok{/}\NormalTok{ denominator,}

    \AttributeTok{half\_width =}\NormalTok{ (}
\NormalTok{      z }\SpecialCharTok{*} \FunctionTok{sqrt}\NormalTok{(}
\NormalTok{        proportion }\SpecialCharTok{*}\NormalTok{ (}\DecValTok{1} \SpecialCharTok{{-}}\NormalTok{ proportion) }\SpecialCharTok{/}\NormalTok{ records }\SpecialCharTok{+}
\NormalTok{          z}\SpecialCharTok{\^{}}\DecValTok{2} \SpecialCharTok{/}\NormalTok{ (}\DecValTok{4} \SpecialCharTok{*}\NormalTok{ records}\SpecialCharTok{\^{}}\DecValTok{2}\NormalTok{)}
\NormalTok{      )}
\NormalTok{    ) }\SpecialCharTok{/}\NormalTok{ denominator,}

    \AttributeTok{lower =}\NormalTok{ centre }\SpecialCharTok{{-}}\NormalTok{ half\_width,}
    \AttributeTok{upper =}\NormalTok{ centre }\SpecialCharTok{+}\NormalTok{ half\_width}
\NormalTok{  )}

\NormalTok{primary\_support }\SpecialCharTok{|\textgreater{}}
  \FunctionTok{ggplot}\NormalTok{(}
    \FunctionTok{aes}\NormalTok{(}
      \AttributeTok{x =}\NormalTok{ dmt\_status,}
      \AttributeTok{y =}\NormalTok{ proportion}
\NormalTok{    )}
\NormalTok{  ) }\SpecialCharTok{+}
  \FunctionTok{geom\_errorbar}\NormalTok{(}
    \FunctionTok{aes}\NormalTok{(}
      \AttributeTok{ymin =}\NormalTok{ lower,}
      \AttributeTok{ymax =}\NormalTok{ upper}
\NormalTok{    ),}
    \AttributeTok{width =} \FloatTok{0.15}\NormalTok{,}
    \AttributeTok{colour =} \StringTok{"\#35193E"}\NormalTok{,}
    \AttributeTok{linewidth =} \FloatTok{0.8}
\NormalTok{  ) }\SpecialCharTok{+}
  \FunctionTok{geom\_point}\NormalTok{(}
    \AttributeTok{colour =} \StringTok{"\#AD1759"}\NormalTok{,}
    \AttributeTok{size =} \FloatTok{3.5}
\NormalTok{  ) }\SpecialCharTok{+}
  \FunctionTok{coord\_flip}\NormalTok{() }\SpecialCharTok{+}
  \FunctionTok{scale\_y\_continuous}\NormalTok{(}
    \AttributeTok{labels =}\NormalTok{ scales}\SpecialCharTok{::}\FunctionTok{label\_percent}\NormalTok{(),}
    \AttributeTok{limits =} \FunctionTok{c}\NormalTok{(}\DecValTok{0}\NormalTok{, }\FloatTok{0.6}\NormalTok{)}
\NormalTok{  ) }\SpecialCharTok{+}
  \FunctionTok{labs}\NormalTok{(}
    \AttributeTok{title =} \StringTok{"Symptom proportions by DMT status"}\NormalTok{,}
    \AttributeTok{subtitle =} \StringTok{"Unadjusted estimates with 95\% Wilson intervals"}\NormalTok{,}
    \AttributeTok{x =} \ConstantTok{NULL}\NormalTok{,}
    \AttributeTok{y =} \StringTok{"Records reporting symptoms"}
\NormalTok{  )}
\end{Highlighting}
\end{Shaded}

\begin{figure}

{\centering \includegraphics[width=1\linewidth]{report_files/figure-latex/exposure-support-1} 

}

\caption{Observed symptom proportions by DMT status with 95\% Wilson confidence intervals.}\label{fig:exposure-support}
\end{figure}

\begin{longtable}[]{@{}lrrrr@{}}
\caption{Outcome support in the secondary analytical
cohort.}\tabularnewline
\toprule\noalign{}
Exposure & Records & Symptoms & No symptoms & Symptoms (\%) \\
\midrule\noalign{}
\endfirsthead
\toprule\noalign{}
Exposure & Records & Symptoms & No symptoms & Symptoms (\%) \\
\midrule\noalign{}
\endhead
\bottomrule\noalign{}
\endlastfoot
Other identified DMT & 588 & 187 & 401 & 31.8 \\
Anti-CD20 & 99 & 33 & 66 & 33.3 \\
\end{longtable}

Both primary and secondary treatment specifications retained variation
in the binary outcome across exposure categories. The secondary cohort
included 33 symptomatic and 66 asymptomatic anti-CD20 records, providing
initial support for coefficient estimation. Nevertheless, unequal
exposure-group sizes and potential sparsity within covariate
combinations required further assessment of model estimability.

\section{Discussion}\label{discussion}

\subsection{Interpretation and fitness for
purpose}\label{interpretation-and-fitness-for-purpose}

\subsection{Limitations and future
research}\label{limitations-and-future-research}

\subsection{Ethics and data
governance}\label{ethics-and-data-governance}

\section{Reproducibility and
Metadata}\label{reproducibility-and-metadata}

\section{Conclusion}\label{conclusion}

\section*{References}\label{references}
\addcontentsline{toc}{section}{References}

\protect\phantomsection\label{refs}
\begin{CSLReferences}{1}{1}
\bibitem[\citeproctext]{ref-evangelou2021}
{Evangelou, Nikos et al.} 2021. {`Self-Diagnosed COVID-19 in People with
Multiple Sclerosis: A Community-Based Cohort of the UK MS Register'}.
\emph{Journal of Neurology, Neurosurgery and Psychiatry}, ahead of
print. \url{https://doi.org/10.1136/jnnp-2020-324449}.

\bibitem[\citeproctext]{ref-khan2024}
Khan, Hamza, Lotte Geys, Peer Baneke, Giancarlo Comi, and Liesbet M.
Peeters. 2024. {`Patient Level Dataset to Study the Effect of COVID-19
in People with Multiple Sclerosis'}. \emph{Scientific Data}, ahead of
print. \url{https://doi.org/10.1038/s41597-024-02978-x}.

\bibitem[\citeproctext]{ref-peeters2020}
{Peeters, L. M. et al.} 2020. {`{COVID-19} in People with Multiple
Sclerosis: A Global Data Sharing Initiative'}. \emph{Multiple Sclerosis
Journal} 26 (10): 1157--62.
\url{https://doi.org/10.1177/1352458520941485}.

\bibitem[\citeproctext]{ref-simpson2022}
{Simpson-Yap, Steve et al.} 2022. {`Updated Results of the COVID-19 in
MS Global Data Sharing Initiative: Anti-CD20 and Other Risk Factors
Associated with COVID-19 Severity'}. \emph{Neurology: Neuroimmunology
and Neuroinflammation}, ahead of print.
\url{https://doi.org/10.1212/NXI.0000000000200021}.

\bibitem[\citeproctext]{ref-sormani2021}
{Sormani, Maria Pia et al.} 2021. {`Disease-Modifying Therapies and
Coronavirus Disease 2019 Severity in Multiple Sclerosis'}. \emph{Annals
of Neurology}, ahead of print. \url{https://doi.org/10.1002/ana.26028}.

\bibitem[\citeproctext]{ref-struyf2022}
{Struyf, Thomas et al.} 2022. {`Signs and Symptoms to Determine If a
Patient Presenting in Primary Care or Hospital Outpatient Settings Has
COVID-19'}. \emph{Cochrane Database of Systematic Reviews}, ahead of
print. \url{https://doi.org/10.1002/14651858.CD013665.pub3}.

\end{CSLReferences}

\end{document}
